\subsubsection{ESTRATÉGIA DE SINCRONIZAÇÃO: \en{CRDT}}

Sincronização de dados em sistemas é uma tarefa essencial no estabelecimento de uma equivalência entre duas ou mais coleções de dados. Do mesmo modo, Tanenbaum \cite{tanenbaum2007dist} afirma que o principal desafio na replicação de dados (isso é, manter diferentes cópias de dados em servidores ou localidades diferentes) é assegurar a consistência das réplicas. 

Estratégias de sincronização estão presentes tanto em arquiteturas centralizadas e descentralizadas. Inicialmente, há os \en{CRDTs} (tipos de dado livre de conflito): algoritmos distribuídos que computam o estado de uma aplicação colaborativa do seu histórico de operações. Esse estado deve ser o mesmo para todos os usuários e ser computável sem a necessidade de ter um componente que centraliza as informações; um servidor central, por exemplo. Possui dois tipos distintos dependendo no que se baseiam: operações ou estados \cite{weidner2023crdtsurvey}.

\en{CRDTs} é um sistema replicado, e como tal precisa propagar atualizações para todas as replicas, há dois caminhos para alcançar esse objetivo: replicação baseada em estado e replicação baseada em operação. No primeiro, a estrutura de dados faz uso de uma função de união que integra o estado da replica remota na local; na segunda forma de replicação, a convergência dos estados ocorre através da propagação de operações entre replicas, portanto uma \en{CRDT} dessa categoria deve gerar, para cada operação, um gerador e um efetor; o primeiro responsável por gerar o segundo ao fim de uma operação; este segundo é, então, executado nas outras replicas para convergir os estados \cite{PreguicaBaqueroShapiro2018CRDT}.

\en{CRDTs} podem ser entendidas como resultado da evolução, até o momento, de estratégias de sincronização em sistemas distribuídos, a aplicação usando o \en{CRDT} de operação é o mais apropriado ao sistema de \en{Offline-First} produzido devido a capacidade das alterações em anotações serem traduzidas em operações (inserir e deletar).

Uma outra forma de sincronização que também faz uso do compartilhamento de estado é a transformação operacional (\en{OT}). Esta é voltada para arquiteturas centralizadas e é muito usada em editores em grupo; ela resolve problemas similares a \en{CRDTs}: consistência na manutenção de documentos compartilhados com tempo de resposta curto e edição paralela em ambientes distribuídos, por exemplo \cite{SunEllis1998OperationalTransformation}. Essa estratégia de sincronização é pesquisada desde o fim dos anos 80 e sua importância é inegável, desde artigos, videos, criação de jogos, comunicação \en{online} em geral. Os desenvolvimentos teóricos e provas experimentais de Ellis et al. que a originaram; atualmente, junto a \en{CRDT} é um dos métodos mais bem aceitos para resolver o problema de consistência de dados em arquiteturas distribuídas \cite{GadeaIonescuIonescu2020ControlLoopOT}.
